\documentclass[11pt,a4paper,openany]{memoir}
\newcommand{\olpath}{../upstream}
\usepackage{fontspec}
\ifLuaTeX
\setmainfont{Nirmala UI}[Script=Tamil,Renderer=HarfBuzz,WordSpace=1.3,ItalicFont={Nirmala UI},ItalicFeatures={FakeSlant=0.15}]
\setsansfont{Nirmala UI}[Script=Tamil,Renderer=HarfBuzz,WordSpace=1.3]
\else
\setmainfont{Nirmala UI}[Script=Tamil,Renderer=OpenType,WordSpace=1.3,ItalicFont={Nirmala UI},ItalicFeatures={FakeSlant=0.15}]
\setsansfont{Nirmala UI}[Script=Tamil,Renderer=OpenType,WordSpace=1.3]
\fi
\setmonofont{Consolas}
\input{\olpath/sty/open-logic.sty}
\ifXeTeX
  \XeTeXlinebreaklocale "ta"
  \XeTeXlinebreakskip=0pt plus 1pt
  \XeTeXgenerateactualtext=1
\fi
\setlrmarginsandblock{25mm}{25mm}{*}
\setulmarginsandblock{22mm}{25mm}{*}
\checkandfixthelayout
\linespread{1.17}
\emergencystretch=2em
% Tamil noun inflection retains the upstream token key and plural/article flags.
% Optional case follows the token: !!{element}[um], !!{element}s[pred].
% Nirmala UI lacks the combining double-acute used by K\H{o}nig.
% Keep the source accent command and render that glyph with a Latin fallback.
\newfontfamily\ollatinaccentfont{Times New Roman}
\let\oloriginaldoubleacute\H
\RenewDocumentCommand\H{m}{{\ollatinaccentfont\oloriginaldoubleacute{#1}}}
\settexttoken{element}{உறுப்பு}{உறுப்புகள்}[உறுப்பு][உறுப்புகள்]
\settexttoken{formula}{வாய்பாடு}{வாய்பாடுகள்}[வாய்பாடு][வாய்பாடுகள்]
\definetoken{p-gen}{formula}{வாய்பாடுகளின்}
\settexttoken{derivation}{வருவித்தல்}{வருவித்தல்கள்}[வருவித்தல்][வருவித்தல்கள்]
\settexttoken{derivability}{வருவிக்கத்தக்கதன்மை}{வருவிக்கத்தக்கதன்மை}[வருவிக்கத்தக்கதன்மை][வருவிக்கத்தக்கதன்மை]
\settexttoken{nonderivability}{வருவிக்கவியலாமை}{வருவிக்கவியலாமை}[வருவிக்கவியலாமை][வருவிக்கவியலாமை]
\settexttoken{sentence}{வாக்கியம்}{வாக்கியங்கள்}[வாக்கியம்][வாக்கியங்கள்]
\definetoken{s-acc}{sentence}{வாக்கியத்தை}
\definetoken{p-acc}{sentence}{வாக்கியங்களை}
\definetoken{s-accum}{sentence}{வாக்கியத்தையும்}
\definetoken{p-accum}{sentence}{வாக்கியங்களையும்}
\definetoken{s-gen}{sentence}{வாக்கியத்தின்}
\definetoken{p-gen}{sentence}{வாக்கியங்களின்}
\definetoken{s-dat}{sentence}{வாக்கியத்திற்கு}
\definetoken{p-dat}{sentence}{வாக்கியங்களுக்கு}
\definetoken{s-loc}{sentence}{வாக்கியத்தில்}
\definetoken{p-loc}{sentence}{வாக்கியங்களில்}
\definetoken{s-abl}{sentence}{வாக்கியத்திலிருந்து}
\definetoken{p-abl}{sentence}{வாக்கியங்களிலிருந்து}
\definetoken{s-ins}{sentence}{வாக்கியத்தால்}
\definetoken{p-ins}{sentence}{வாக்கியங்களால்}
\definetoken{s-com}{sentence}{வாக்கியத்துடன்}
\definetoken{p-com}{sentence}{வாக்கியங்களுடன்}
\definetoken{s-acc}{derivation}{வருவித்தலை}
\definetoken{p-acc}{derivation}{வருவித்தல்களை}
\definetoken{s-gen}{derivation}{வருவித்தலின்}
\definetoken{p-gen}{derivation}{வருவித்தல்களின்}
\definetoken{s-loc}{derivation}{வருவித்தலில்}
\definetoken{p-loc}{derivation}{வருவித்தல்களில்}
\settexttoken{derive}{வருவிக்க}{வருவிக்க}[வருவிக்க][வருவிக்க]
\settexttoken{derivable}{வருவிக்கத்தக்க}{வருவிக்கத்தக்க}[வருவிக்கத்தக்க][வருவிக்கத்தக்க]
\settexttoken{tableau}{அட்டவணை மரம்}{அட்டவணை மரங்கள்}[அட்டவணை மரம்][அட்டவணை மரங்கள்]
\definetoken{s-gen}{tableau}{அட்டவணை மரத்தின்}
\settexttoken{signed formula}{குறியிட்ட வாய்பாடு}{குறியிட்ட வாய்பாடுகள்}[குறியிட்ட வாய்பாடு][குறியிட்ட வாய்பாடுகள்]
\settexttoken{discharge}{விடுவித்தல்}{விடுவித்தல்கள்}[விடுவித்தல்][விடுவித்தல்கள்]
\settexttoken{discharged}{விடுவிக்கப்பட்ட}{விடுவிக்கப்பட்ட}[விடுவிக்கப்பட்ட][விடுவிக்கப்பட்ட]
\settexttoken{undischarged}{விடுவிக்கப்படாத}{விடுவிக்கப்படாத}[விடுவிக்கப்படாத][விடுவிக்கப்படாத]
\settexttoken{variable}{மாறி}{மாறிகள்}[மாறி][மாறிகள்]
\definetoken{p-dat}{variable}{மாறிகளுக்குப்}
\settexttoken{constant}{மாறிலிக் குறி}{மாறிலிக் குறிகள்}[மாறிலிக் குறி][மாறிலிக் குறிகள்]
\settexttoken{predicate}{பயனிலைக் குறி}{பயனிலைக் குறிகள்}[பயனிலைக் குறி][பயனிலைக் குறிகள்]
\settexttoken{function}{சார்புக் குறி}{சார்புக் குறிகள்}[சார்புக் குறி][சார்புக் குறிகள்]
\settexttoken{domain}{பொருட்களம்}{பொருட்களங்கள்}[பொருட்களம்][பொருட்களங்கள்]
\settexttoken{value}{மதிப்பு}{மதிப்புகள்}[மதிப்பு][மதிப்புகள்]
\settexttoken{main operator}{முதன்மைச் செயலி}{முதன்மைச் செயலிகள்}[முதன்மைச் செயலி][முதன்மைச் செயலிகள்]
% A term is "free for" a variable when substituting it causes no capture.
% Keep this as a predicate phrase so protected !!{free for} tokens read
% naturally after both singular and plural mathematical subjects.
\settexttoken{free for}{இடைநிறுத்தத்தக்கது}{இடைநிறுத்தத்தக்கவை}[இடைநிறுத்தத்தக்கது][இடைநிறுத்தத்தக்கவை]
\settexttoken{complete}{நிறைவான}{நிறைவான}[நிறைவான][நிறைவான]
\settexttoken{propositional variable}{கூற்று மாறி}{கூற்று மாறிகள்}[கூற்று மாறி][கூற்று மாறிகள்]
\settexttoken{valuation}{மெய்மதிப்பு ஒதுக்கீடு}{மெய்மதிப்பு ஒதுக்கீடுகள்}[மெய்மதிப்பு ஒதுக்கீடு][மெய்மதிப்பு ஒதுக்கீடுகள்]
\settexttoken{conditional}{நிபந்தனைக் கூற்று}{நிபந்தனைக் கூற்றுகள்}[நிபந்தனைக் கூற்று][நிபந்தனைக் கூற்றுகள்]
\settexttoken{biconditional}{இரு நிபந்தனைக் கூற்று}{இரு நிபந்தனைக் கூற்றுகள்}[இரு நிபந்தனைக் கூற்று][இரு நிபந்தனைக் கூற்றுகள்]
\settexttoken{falsity}{பொய்மை}{பொய்மைகள்}[பொய்மை][பொய்மைகள்]
\settexttoken{truth}{மெய்மை}{மெய்மைகள்}[மெய்மை][மெய்மைகள்]
\settexttoken{operator}{செயலி}{செயலிகள்}[செயலி][செயலிகள்]
\definetoken{p-dat}{operator}{செயலிகளுக்குச்}
\settexttoken{structure}{கட்டமைப்பு}{கட்டமைப்புகள்}[கட்டமைப்பு][கட்டமைப்புகள்]
\definetoken{p-loc}{structure}{கட்டமைப்புகளில்}
\settexttoken{identity}{ஒன்றாதல்}{ஒன்றாதல்கள்}[ஒன்றாதல்][ஒன்றாதல்கள்]
\settexttoken{denumerable}{எண்ணிடத்தக்க}{எண்ணிடத்தக்க}[எண்ணிடத்தக்க][எண்ணிடத்தக்க]
\definetoken{a}{denumerable}{ஓர்}
\definetoken{A}{denumerable}{ஓர்}
\definetoken{a}{propositional variable}{ஒரு}
\definetoken{A}{propositional variable}{ஒரு}
\definetoken{a}{valuation}{ஒரு}
\definetoken{A}{valuation}{ஒரு}
\definetoken{a}{conditional}{ஒரு}
\definetoken{A}{conditional}{ஒரு}
\definetoken{a}{biconditional}{ஓர்}
\definetoken{A}{biconditional}{ஓர்}
\definetoken{a}{structure}{ஒரு}
\definetoken{A}{structure}{ஒரு}
\definetoken{a}{formula}{ஒரு}
\definetoken{A}{formula}{ஒரு}
\definetoken{a}{derivation}{ஒரு}
\definetoken{A}{derivation}{ஒரு}
\definetoken{a}{sentence}{ஒரு}
\definetoken{A}{sentence}{ஒரு}
\definetoken{a}{tableau}{ஓர்}
\definetoken{A}{tableau}{ஓர்}
\definetoken{a}{signed formula}{ஒரு}
\definetoken{A}{signed formula}{ஒரு}
\definetoken{a}{variable}{ஒரு}
\definetoken{A}{variable}{ஒரு}
\definetoken{a}{constant}{ஒரு}
\definetoken{A}{constant}{ஒரு}
\definetoken{a}{predicate}{ஒரு}
\definetoken{A}{predicate}{ஒரு}
\definetoken{a}{function}{ஒரு}
\definetoken{A}{function}{ஒரு}
\definetoken{a}{domain}{ஒரு}
\definetoken{A}{domain}{ஒரு}
\definetoken{a}{value}{ஒரு}
\definetoken{A}{value}{ஒரு}
\definetoken{a}{complete}{ஒரு}
\definetoken{A}{complete}{ஒரு}
\settexttoken{injective}{ஒன்றுக்கொன்றான}{ஒன்றுக்கொன்றான}[ஒன்றுக்கொன்றான][ஒன்றுக்கொன்றான]
\settexttoken{surjective}{மேற்கோர்த்தல்}{மேற்கோர்த்தல்}[மேற்கோர்த்தல்][மேற்கோர்த்தல்]
\settexttoken{enumerable}{எண்ணிடத்தக்க}{எண்ணிடத்தக்க}[எண்ணிடத்தக்க][எண்ணிடத்தக்க]
\settexttoken{nonenumerable}{எண்ணிடத்தகாத}{எண்ணிடத்தகாத}[எண்ணிடத்தகாத][எண்ணிடத்தகாத]
\definetoken{a}{enumerable}{ஓர்}
\definetoken{A}{enumerable}{ஓர்}
\definetoken{a}{nonenumerable}{ஓர்}
\definetoken{A}{nonenumerable}{ஓர்}
\settexttoken{bijective}{இருபுற}{இருபுற}[இருபுற][இருபுற]
\settexttoken{injection}{ஒன்றுக்கொன்றான சார்பு}{ஒன்றுக்கொன்றான சார்புகள்}[ஒன்றுக்கொன்றான சார்பு][ஒன்றுக்கொன்றான சார்புகள்]
\settexttoken{surjection}{மேற்கோர்த்தல் சார்பு}{மேற்கோர்த்தல் சார்புகள்}[மேற்கோர்த்தல் சார்பு][மேற்கோர்த்தல் சார்புகள்]
\settexttoken{bijection}{இருபுறச் சார்பு}{இருபுறச் சார்புகள்}[இருபுறச் சார்பு][இருபுறச் சார்புகள்]
\settexttoken{computably enumerable}{கணிக்கத்தக்கவாறு எண்ணிடத்தக்க}{கணிக்கத்தக்கவாறு எண்ணிடத்தக்க}[கணிக்கத்தக்கவாறு எண்ணிடத்தக்க][கணிக்கத்தக்கவாறு எண்ணிடத்தக்க]
\settexttoken{c.e.}{க.வ.எ.}{க.வ.எ.}[க.வ.எ.][க.வ.எ.]
\settexttoken{decidable}{தீர்மானிக்கத்தக்க}{தீர்மானிக்கத்தக்க}[தீர்மானிக்கத்தக்க][தீர்மானிக்கத்தக்க]
\settexttoken{axiomatizable}{அடிகோளாக்கத்தக்க}{அடிகோளாக்கத்தக்க}[அடிகோளாக்கத்தக்க][அடிகோளாக்கத்தக்க]
\settexttoken{axiomatized}{அடிகோளாக்கப்பட்ட}{அடிகோளாக்கப்பட்ட}[அடிகோளாக்கப்பட்ட][அடிகோளாக்கப்பட்ட]
\settexttoken{axiomatizability}{அடிகோளாக்கத்தன்மை}{அடிகோளாக்கத்தன்மை}[அடிகோளாக்கத்தன்மை][அடிகோளாக்கத்தன்மை]
\settexttoken{language}{மொழி}{மொழிகள்}[மொழி][மொழிகள்]
\settexttoken{represents}{பிரதிநிதித்துவப்படுத்துகிறது}{பிரதிநிதித்துவப்படுத்துகின்றன}[பிரதிநிதித்துவப்படுத்துகிறது][பிரதிநிதித்துவப்படுத்துகின்றன]
\settexttoken{subformula}{உட்வாய்பாடு}{உட்வாய்பாடுகள்}[உட்வாய்பாடு][உட்வாய்பாடுகள்]
\settexttoken{lambda define}{லாம்டா-வரையறுக்கிறது}{லாம்டா-வரையறுக்கின்றன}[லாம்டா-வரையறுக்கிறது][லாம்டா-வரையறுக்கின்றன]
\settexttoken{lambda defined}{லாம்டா-வரையறுக்கப்பட்ட}{லாம்டா-வரையறுக்கப்பட்ட}[லாம்டா-வரையறுக்கப்பட்ட]
\settexttoken{lambda definable}{லாம்டா-வரையறுக்கத்தக்க}{லாம்டா-வரையறுக்கத்தக்க}[லாம்டா-வரையறுக்கத்தக்க]
\definetoken{a}{language}{ஒரு}
\definetoken{A}{language}{ஒரு}
\definetoken{a}{subformula}{ஓர்}
\definetoken{A}{subformula}{ஓர்}
\definetoken{a}{injective}{ஓர்}
\definetoken{a}{surjective}{ஒரு}
\definetoken{a}{bijective}{ஓர்}
\definetoken{a}{injection}{ஓர்}
\definetoken{a}{surjection}{ஒரு}
\definetoken{a}{bijection}{ஓர்}
\definetoken{A}{injective}{ஓர்}
\definetoken{A}{surjective}{ஒரு}
\definetoken{A}{bijective}{ஓர்}
\definetoken{A}{injection}{ஓர்}
\definetoken{A}{surjection}{ஒரு}
\definetoken{A}{bijection}{ஓர்}
\definetoken{a}{element}{ஓர்}
\definetoken{A}{element}{ஓர்}
\definetoken{s-um}{element}{உறுப்பும்}
\definetoken{p-um}{element}{உறுப்புகளும்}
\definetoken{s-pred}{element}{உறுப்பாக}
\definetoken{p-pred}{element}{உறுப்புகளாக}
\definetoken{s-acc}{element}{உறுப்பை}
\definetoken{p-acc}{element}{உறுப்புகளை}
\definetoken{s-gen}{element}{உறுப்பின்}
\definetoken{p-gen}{element}{உறுப்புகளின்}
\definetoken{s-loc}{element}{உறுப்பில்}
\definetoken{p-loc}{element}{உறுப்புகளில்}
% Cumulative-hierarchy principle names used throughout the set-theory part.
\renewcommand\stageshier{\emph{கணங்கள் படிநிலைகளில் தோன்றுகின்றன}}
\renewcommand\stagesord{\emph{படிநிலைகள் வரிசைப்படுத்தப்படுகின்றன}}
\renewcommand\stagesacc{\emph{படிநிலைகள் திரள்கின்றன}}
\renewcommand\stagessucc{\emph{படிநிலைகள் தொடர்கின்றன}}
\renewcommand\stagesinf{\emph{படிநிலைகள் முடிவிலியை அடைகின்றன}}
\makeatletter
\RenewDocumentCommand \@printtoken { t^ ta m !ts O{} }{%
  \begingroup
  \IfBooleanTF{#4}{\def\TAForm{p}}{\def\TAForm{s}}%
  \ifstrempty{#5}{}{\edef\TAForm{\TAForm-#5}}%
  \IfBooleanT{#2}{\article{#3}~}%
  \printtoken{\TAForm}{#3}%
  \endgroup}
\makeatother

% Preserve extraction of composite negated relation glyphs without changing source formulas.
\usepackage{accsupp}
\let\TAOriginalNotin\notin
\let\TAOriginalNeq\neq
\let\TAOriginalNsubseteq\nsubseteq
\renewcommand{\notin}{\mathrel{\BeginAccSupp{method=hex,unicode,ActualText=2209}\TAOriginalNotin\EndAccSupp{}}}
\renewcommand{\neq}{\mathrel{\BeginAccSupp{method=hex,unicode,ActualText=2260}\TAOriginalNeq\EndAccSupp{}}}
\renewcommand{\nsubseteq}{\mathrel{\BeginAccSupp{method=hex,unicode,ActualText=2288}\TAOriginalNsubseteq\EndAccSupp{}}}

\tagtrue{novice,math,compsci,FOL,prvTrue,prvFalse,prvEx,prvAll}
\addto\captionsenglish{%
  \renewcommand{\chaptername}{அத்தியாயம்}%
  \renewcommand{\partname}{பகுதி}%
  \renewcommand{\appendixname}{இணைப்பு}%
  \renewcommand{\tablename}{அட்டவணை}%
}
\setlocalecaption{english}{photocredits}{படங்களுக்கான நன்றியுரை}
\setlocalecaption{english}{history}{வரலாற்றுக் குறிப்புகள்}
\setlocalecaption{english}{reading}{மேலும் படிக்க}
\setlocalecaption{english}{appendix}{இணைப்பு}
\setlocalecaption{english}{appendices}{இணைப்புகள்}
\setlocalecaption{english}{definition}{வரையறை}
\setlocalecaption{english}{example}{எடுத்துக்காட்டு}
\setlocalecaption{english}{lemma}{துணைத்தேற்றம்}
\setlocalecaption{english}{proposition}{கூற்று}
\setlocalecaption{english}{corollary}{தொடர்விளைவு}
\setlocalecaption{english}{problem}{பயிற்சி}
\setlocalecaption{english}{problems}{பயிற்சிகள்}
\setlocalecaption{english}{theorem}{தேற்றம்}
\setlocalecaption{english}{remark}{குறிப்பு}
\setlocalecaption{english}{axiom}{அடிக்கோள்}
\setlocalecaption{english}{note}{குறிப்பு}
\setlocalecaption{english}{case}{நிலை}
\setlocalecaption{english}{convention}{மரபு}
\addto\captionsenglish{\renewcommand{\proofname}{நிறுவல்}}
\addto\captionsenglish{\renewcommand{\figurename}{படம்}}
\crefname{page}{பக்கம்}{பக்கங்கள்}
\Crefname{page}{பக்கம்}{பக்கங்கள்}
\crefname{part}{பகுதி}{பகுதிகள்}
\Crefname{part}{பகுதி}{பகுதிகள்}
\crefname{chapter}{அத்தியாயம்}{அத்தியாயங்கள்}
\Crefname{chapter}{அத்தியாயம்}{அத்தியாயங்கள்}
\crefname{section}{பிரிவு}{பிரிவுகள்}
\Crefname{section}{பிரிவு}{பிரிவுகள்}
\crefname{subsection}{பிரிவு}{பிரிவுகள்}
\Crefname{subsection}{பிரிவு}{பிரிவுகள்}
\crefname{figure}{படம்}{படங்கள்}
\Crefname{figure}{படம்}{படங்கள்}
\crefname{table}{அட்டவணை}{அட்டவணைகள்}
\Crefname{table}{அட்டவணை}{அட்டவணைகள்}
\crefname{equation}{சமன்பாடு}{சமன்பாடுகள்}
\Crefname{equation}{சமன்பாடு}{சமன்பாடுகள்}
\crefname{defn}{வரையறை}{வரையறைகள்}
\Crefname{defn}{வரையறை}{வரையறைகள்}
\crefname{thm}{தேற்றம்}{தேற்றங்கள்}
\Crefname{thm}{தேற்றம்}{தேற்றங்கள்}
\crefname{ex}{எடுத்துக்காட்டு}{எடுத்துக்காட்டுகள்}
\Crefname{ex}{எடுத்துக்காட்டு}{எடுத்துக்காட்டுகள்}
\crefname{lem}{துணைத்தேற்றம்}{துணைத்தேற்றங்கள்}
\Crefname{lem}{துணைத்தேற்றம்}{துணைத்தேற்றங்கள்}
\crefname{prop}{கூற்று}{கூற்றுகள்}
\Crefname{prop}{கூற்று}{கூற்றுகள்}
\crefname{prob}{பயிற்சி}{பயிற்சிகள்}
\Crefname{prob}{பயிற்சி}{பயிற்சிகள்}
\crefname{rem}{குறிப்பு}{குறிப்புகள்}
\Crefname{rem}{குறிப்பு}{குறிப்புகள்}
\crefname{cor}{தொடர்விளைவு}{தொடர்விளைவுகள்}
\Crefname{cor}{தொடர்விளைவு}{தொடர்விளைவுகள்}
\crefname{axiom}{அடிக்கோள்}{அடிக்கோள்கள்}
\Crefname{axiom}{அடிக்கோள்}{அடிக்கோள்கள்}
\crefname{note}{குறிப்பு}{குறிப்புகள்}
\Crefname{note}{குறிப்பு}{குறிப்புகள்}
\crefname{case}{நிலை}{நிலைகள்}
\Crefname{case}{நிலை}{நிலைகள்}
\crefname{conv}{மரபு}{மரபுகள்}
\Crefname{conv}{மரபு}{மரபுகள்}
\AtBeginEnvironment{align*}{\fontsize{9.5}{12}\selectfont}
\hypersetup{pdftitle={OpenLogic Tamil: sets through sequent calculus},pdfauthor={OpenLogic Tamil translation programme},pdfsubject={Machine translation; 80 of 722 frozen source units}}
\addto\captionsenglish{\renewcommand{\contentsname}{பொருளடக்கம்}\renewcommand{\bibname}{மேற்கோள் நூல்}}
\begin{document}
\raggedright
\pagestyle{plain}
\chapter*{திறந்த தருக்கவியல்}
\noindent தமிழ் மொழிபெயர்ப்பு — கணங்களிலிருந்து தொடரணிக் கணியம் வரை

\noindent மூல ஆசிரியர்கள்: Open Logic Project.
\href{https://github.com/OpenLogicProject/OpenLogic}{மூல நூல்},
\href{https://creativecommons.org/licenses/by/4.0/}{CC BY 4.0}.
இந்தப் பிரதி இயந்திர மொழிபெயர்ப்பு; சுயாதீன மனிதச் சரிபார்ப்பு செய்யப்பட்டதாகக்
கூறப்படவில்லை. முழுப் பதிப்பு இன்னும் தயாராகவில்லை.

\noindent இந்தப் பிரதியில் 722 உறையவைக்கப்பட்ட மூல அலகுகளில் 80
மொழிபெயர்க்கப்பட்டுள்ளன. கணங்கள், தொடர்புகள், சார்புகள், கணங்களின் அளவு,
எண்முறைமைகளின் கணக்கோட்பாட்டுக் கட்டமைப்பு, முடிவுறாத கணங்கள் ஆகியவற்றுடன்,
கூற்றுத் தருக்கவியலின் தொடரமைப்பும் பொருண்மையியலும் இடம்பெறுகின்றன. நிறுவல்
முறைமைகள் பற்றிய மேலோட்டத்தையும், விதிகள், நிறுவல் தேடல், அளவையடைகள்,
முரணின்மை, சரித்தன்மை, ஒன்றாதல் ஆகியவற்றுடன் செவ்வியல் முதல்-வரிசைத்
தருக்கவியலுக்கான முழுத் தொடரணிக் கணிய அதிகாரத்தையும் இப்பிரதி கொண்டுள்ளது.

\clearpage
\tableofcontents*
\subfile{../translation/content/sets-functions-relations/sets/sets.tex}
\section*{தமிழ்ப் பதிப்பின் குறிப்பு}
இந்தக் குறிப்பு மூல நூலின் பகுதியல்ல.
இந்நூல் இயல் எண்களின் கணத்தில் பூச்சியத்தையும் சேர்க்கும் மரபைப்
பின்பற்றுகிறது: $\Nat=\{0,1,2,\ldots\}$. தமிழ்நாடு ஆறாம் வகுப்புக்
கணக்குப் பாடநூலின் 2018 முதல் பருவப் பதிப்பில், அச்சுப் பக்கம் 29 இல்,
இயல் எண்கள் 1 இலிருந்து தொடங்குகின்றன; பூச்சியத்தையும் சேர்த்த கணம்
முழு எண்களின் கணம் எனப்படுகிறது. இங்கே OpenLogic மூல நூலின்
மரபே பின்பற்றப்படுகிறது. எதிர்ம எண்களையும் உள்ளடக்கிய
$\Int=\{\ldots,-1,0,1,\ldots\}$ இன் உறுப்புகளுக்கு, அதே பள்ளி நூலின்
அச்சுப் பக்கம் 169 இல் தரப்பட்டுள்ளபடி, “முழுக்கள்” என்ற சொல்லைப்
பயன்படுத்துகிறோம். “முழு எண்கள்” என்பதுடன் இதைக் குழப்ப வேண்டாம்.

\subfile{../translation/content/sets-functions-relations/relations/relations-complete.tex}
\section*{மூல உரை பற்றிய பதிப்புக் குறிப்புகள்}
இந்தக் குறிப்புகள் மூல நூலின் பகுதிகள் அல்ல. இவை உறுதிசெய்யப்பட்ட
ஆங்கில மூலப் பதிப்பில் உள்ள சில குறியீட்டு நுணுக்கங்களை விளக்குகின்றன.

மரவுருவின் கிளை வரையறையில் $z\in X\setminus B$ என்று உள்ளது.
அங்கு $X$ தனியாக வரையறுக்கப்படவில்லை; முன்பு அறிமுகப்படுத்தப்பட்ட
மரவுருவின் அடிப்படைக் கணம் $A$ ஆகும். மூலக் குறியீடு மாற்றப்படாமல்
வைக்கப்பட்டுள்ளது.

$\Nat^*$ இன் உள்மரவுரு பற்றிய எடுத்துக்காட்டில், தொடக்கத் துண்டுகளின் கீழ்
அடைவு கொண்ட $A$ பற்றிய கூற்றுக்கு, $A$ வெற்றற்றது என்ற நிபந்தனையும்
தேவைப்படுகிறது. ஏனெனில் இந்நூலின் மரவுரு வரையறை ஒரு வேரைக் கோருகிறது.

வரிசைகள் பற்றிய பிரிவில் $R^+$ என்பது தற்சுட்டு அடைவைக் குறிக்கிறது;
தொடர்புகள் மீதான செயல்கள் பற்றிய பிரிவில் அதே குறியீடு கடப்பு அடைவைக்
குறிக்கிறது. அந்தந்த இடத்தில் தரப்பட்ட வரையறையையே பயன்படுத்த வேண்டும்.

\subfile{../translation/content/sets-functions-relations/functions/functions.tex}
\section*{சார்புகள் பற்றிய பதிப்புக் குறிப்புகள்}
இந்தக் குறிப்புகள் மூல நூலின் பகுதிகள் அல்ல.
\textbf{OLFUN-002.} வர்க்கமூலச் சார்பு பற்றிய எடுத்துக்காட்டில்
“நேர்ம வர்க்க மூலம்” என்று உறையவைக்கப்பட்ட ஆங்கில மூலம் கூறுகிறது.
ஆனால் அதன் சார்பகத்தில் பூச்சியமும் உள்ளது; பூச்சியத்தின் வர்க்க மூலம்
பூச்சியம். ஆகவே மொழிபெயர்க்கப்பட்ட உடல் உரை இதனை
“எதிர்மமற்ற முதன்மை வர்க்க மூலம்” எனத் திருத்துகிறது.
காட்சிப்படுத்தப்பட்ட சார்பு மாறவில்லை.

\textbf{OLFUN-003.} $g$ என்ற சார்பின் எடுத்துக்காட்டில்,
தரப்பட்ட இயல் எண்ணை உறையவைக்கப்பட்ட ஆங்கில மூலம் $n$ எனக்
குறிப்பிட்டுவிட்டு, அடுத்த விளக்கத்தில் $x$ எனக் குறிப்பிடுகிறது.
தமிழ் உடல் உரை தேவையற்ற மாறிப் பெயரை நீக்குகிறது.
$g(x)=x+2-1$ என்ற வாய்பாடும் அதன் கணிதப் பொருளும் மாறவில்லை.

\textbf{OLFUN-004.} ஒரு சார்பின் கோட்டுருவை உறையவைக்கப்பட்ட
ஆங்கில மூலம் ஓரிடத்தில் “$A\times B$ மீதான தொடர்பு” என்கிறது.
இந்நூலின் சொந்த வரையறையின்படி அத்தொடர் வேறு பொருளைத் தரும்.
ஆகவே தமிழ் உடல் உரை அதனை “$A$ க்கும் $B$ க்கும் இடையிலான
$R_f\subseteq A\times B$ என்ற தொடர்பு” எனத் திருத்துகிறது.

\textbf{OLFUN-005.} தொடர்புகளுக்கான முந்தைய கட்டுப்படுத்தல் வரையறை
$R\cap C^2$ ஆகும்;
இது ஜோடியின் இரண்டு உறுப்புகளையும் $C$ க்குள் கட்டுப்படுத்துகிறது.
சார்புகளுக்கான கட்டுப்படுத்தல் வரையறை உள்ளீட்டை மட்டும் $C$ க்குள்
கட்டுப்படுத்துகிறது. ஆகவே பொதுவாக இவை ஒரே செயலல்ல:
சார்பின் வெளியீடு $C$ க்கு வெளியேயும் இருக்கலாம்.
சரியான இரு வாய்பாடுகளும் வைக்கப்பட்டுள்ளன; அவை “அப்படியே”
ஒத்தவை என்ற ஆங்கில விளக்கம் தமிழ் உடல் உரையில் துல்லியமாகத்
தகுதிப்படுத்தப்பட்டுள்ளது.

\textbf{OLFUN-001.} ஒன்றுக்கொன்றான சார்புக்கு இடது நேர்மாறு
உள்ளது என்ற உறையவைக்கப்பட்ட ஆங்கில முன்மொழிவின் நிறுவல்,
சார்பகம் $A$ இலிருந்து ஓர் உறுப்பைத் தேர்ந்தெடுக்கிறது.
இதற்கு $A\neq\emptyset$ என்ற நிபந்தனை தேவை; தமிழ் உடல்
முன்மொழிவிலும் நிறுவலிலும் அது சேர்க்கப்பட்டுள்ளது.
$A=\varnothing$, $B\ne\varnothing$ எனில்,
$f\colon A\to B$ என்ற வெற்றுச் சார்பு ஒன்றுக்கொன்றானது;
ஆனால் $B\to A$ என்ற சார்பு எதுவும் இல்லை.
$A=B=\varnothing$ என்ற நிலையில் வெற்றுச் சார்பே அதன் நேர்மாறு.
துல்லியமான பொதுநிபந்தனை: ஒன்றுக்கொன்றான
$f\colon A\to B$ க்கு இடது நேர்மாறு இருப்பதற்குத் தேவையானதும்
போதுமானதுமான நிபந்தனை $A\neq\emptyset$ அல்லது $B=\emptyset$ ஆகும்.
உறையவைக்கப்பட்ட ஆங்கிலக் கோப்பு மாற்றப்படவில்லை.

\subfile{../translation/content/sets-functions-relations/size-of-sets/size-of-sets-complete.tex}
\section*{கணங்களின் அளவு பற்றிய பதிப்புக் குறிப்புகள்}
இந்தக் குறிப்புகள் மூல நூலின் பகுதிகள் அல்ல.

\textbf{TA-SIZ-LOCAL-001.} ஜோடியாக்கச் சார்புகள் பற்றிய பிரிவில்,
$k$ ஆவது முக்கோண எண் “$k$ ஐ விடச் சிறிய இயல் எண்களின் கூட்டுத்தொகை”
என்றும், அதன் வாய்பாடு $k(k+1)/2$ என்றும் உறையவைக்கப்பட்ட ஆங்கில மூலம்
கூறுகிறது. இந்த நூலில் இயல் எண்கள் பூச்சியத்திலிருந்து தொடங்குவதால்,
$k$ ஐ விடச் சிறிய இயல் எண்களின் கூட்டுத்தொகை உண்மையில்
$k(k-1)/2$ ஆகும். ஆனால் அந்தப் பிரிவின் $0,1,3,6,\ldots$ என்ற
வரிசையும் அதன் ஜோடியாக்க வாய்பாடும் $k(k+1)/2$ என்ற
வழக்கமான பூச்சியம்-அடிப்படை முக்கோண எண் குறியிடலைப் பயன்படுத்துகின்றன.
மூல வாக்கியமும் வாய்பாடுகளும் உடல் உரையில் வைக்கப்பட்டுள்ளன;
இது குறியெண் வாய்பாட்டின் மதிப்புகளை மாற்றுவதில்லை.

\textbf{TA-SIZ-LOCAL-002.} மாற்று ஜோடியாக்கத்தின் முதலாவது படியில்,
“இரண்டாவது” ஜோடியை உறையவைக்கப்பட்ட ஆங்கில மூலம்
$\tuple{0,2}$ என்கிறது; உடனுள்ள அட்டவணையில் அது
$\tuple{0,1}$ ஆக உள்ளது. அட்டவணைக்கும் தொடர்வரிசைக்கும் பொருந்துமாறு
தமிழ் உடல் $\tuple{0,1}$ எனத் திருத்துகிறது.

\textbf{OLSIZ-001.} முழுக்களைப் பட்டியலாக்கும் அட்டவணையின் தலைப்பு
$f(7)$ வரை செல்கிறது; அதற்கு ஒத்த வாய்பாடு
$-\lceil 6/2\rceil=-3$. ஆனால் உறையவைக்கப்பட்ட ஆங்கில மூலத்தின்
மதிப்பு வரிசை $3$ உடன் நின்று $-3$ ஐ விடுகிறது.
தமிழ் அட்டவணை $f(7)$ இன் மதிப்பாக $-3$ ஐச் சேர்க்கிறது.

\textbf{OLSIZ-002.} இணைமுடிவுறு உட்கணத்தின் வரையறையில்,
உறையவைக்கப்பட்ட ஆங்கில மூலத்திலுள்ள “ஒரு முடிவுறு கணம் $\Nat$ இன்
நிரப்பி” என்ற சொற்றொடர், முடிவுறு கணத்திற்கும் $\Nat$ க்கும் உள்ள
உறவை விடுகிறது. உடனுள்ள $\Nat\setminus A$ முடிவுறும் என்ற
சமானநிபந்தனைக்கு ஏற்ப, தமிழ் உடல் இதனை $\Nat$ இன் ஒரு முடிவுறு
உட்கணத்திற்கு, இயல் எண்களின் கணத்திற்குள்ளான நிரப்பி எனத் தெளிவாக்குகிறது.

\textbf{OLSIZ-003.} மாற்று ஜோடியாக்கக் கட்டுமானத்தின் அடுத்த படியில்
$\tuple{2,m}$ என்பதை ஆங்கில மூலம் இருமுறை அடுத்தடுத்து அச்சிடுகிறது.
பின்னுள்ள நிறைவடைந்த பட்டியலில் கிடைவரிசைகள் $2$, $3$ என்று
தொடர்வதால், தமிழ் உடல் இரண்டாவதை $\tuple{3,m}$ எனத் திருத்துகிறது.

\textbf{OLSIZ-004.} முதன்மைக் குறைத்தல் நிறுவலில்,
$f(Z)$ இன் வெளியீட்டை $s_k$ என்று உறையவைக்கப்பட்ட ஆங்கில மூலம்
பெயரிடுகிறது; ஆனால் $k$ அங்கு கட்டுப்படுத்தப்படவோ தீர்மானிக்கப்படவோ
இல்லை. அடுத்த உறுப்புவாரி நிபந்தனை $Z$ இன் இயல்புக்குறித் தொடரைத்
தனித்த முறையில் தீர்மானிக்கிறது. தமிழ் உடல் அந்த வெளியீட்டையும் அதன்
உறுப்புகளையும் ஒரே கட்டுப்பட்ட பெயரான $s$, $s(n)$ கொண்டு குறிக்கிறது.

\textbf{OLSIZ-005.} அதே பிரிவின் எதிரெடுத்துக்காட்டில்,
$h\colon\PosInt\to\Bin^\omega$ என அறிவித்த பிறகு, $h(n)$ ஐ
$n$ பூச்சியங்களைக் கொண்ட முடிவுறு தொடராக உறையவைக்கப்பட்ட ஆங்கில மூலம்
காட்டுகிறது. அது $\Bin^\omega$ இன் உறுப்பு அல்ல. தமிழ் உடல்,
$n$ பூச்சியங்களுக்குப் பின் முடிவுறாமல் ஒன்றுகள் வரும் தொடரைப்
பயன்படுத்துகிறது; ஆகவே ஒவ்வொரு வெளியீடும் உண்மையாகவே
$\Bin^\omega$ இல் உள்ளது.

\textbf{OLSIZ-006.} எண்ணொத்த தன்மை பற்றிய நிறுவலில்,
$f\colon A\to B$ என்ற இருபுறச் சார்பை முன்கொண்ட பிறகு,
$A=\emptyset$ எனில் $B=\emptyset$ என்பதை விளக்கும் இரண்டு
நிபந்தனைக் கிளைகளிலும் உறையவைக்கப்பட்ட ஆங்கில மூலம்
$g(x)=y$ என்கிறது. அந்த நிலையில் தேவையான சார்பு $f$ ஆகும்;
$g$ என்பது பின்னர் பயன்படுத்தப்படும் பட்டியலாக்கம்.
தமிழ் உடல் இரண்டு இடங்களிலும் $f(x)=y$ எனத் திருத்துகிறது.

\textbf{OLSIZ-007.} காண்டோரின் தேற்றத்தின் விரிவான நிறுவலில்,
ஏதேனும் ஒரு $x\in A$ க்கு $g(x)\ne\overline A$ எனக்
காட்ட வேண்டிய இடத்தில், உறையவைக்கப்பட்ட ஆங்கில மூலம்
“ஒவ்வொரு $x\in\overline A$ க்கும்” என்கிறது.
முந்தைய ஏதேனும்-ஒன்றான தேர்வுக்கும் அடுத்த முடிவுக்கும் பொருந்துமாறு,
தமிழ் உடல் “ஒவ்வொரு $x\in A$ க்கும்” எனத் திருத்துகிறது.

\textbf{OLSIZ-008.} மாற்று எண்ணிடத்தகாமைப் பிரிவின் மூலைவிட்ட
அட்டவணையை அறிமுகப்படுத்தும்போது, $s_n(m)$ என்பதை உறையவைக்கப்பட்ட
ஆங்கில மூலம் “$m$ ஆவது தொடரின் $n$ ஆவது இலக்கம்” என்கிறது.
ஆனால் உடனுள்ள அட்டவணையில் $s_n$ என்பது $n$ ஆவது கிடைவரிசைத் தொடரும்,
$m$ என்பது அதன் நெடுவரிசை இலக்கமும் ஆகும்.
தமிழ் உடல் இதனை “$n$ ஆவது தொடரின் $m$ ஆவது இலக்கம்” எனத் திருத்துகிறது.

\textbf{OLSIZ-009.} அதே நிறுவலின் மூலைவிட்டக் கட்டுமானத்தில்,
ஒவ்வொரு $1$ ஐயும் $0$ ஆகவும் ஒவ்வொரு $1$ ஐயும் $0$ ஆகவும்
மாற்றுவதாக ஆங்கில மூலம் இருமுறை ஒரே மாற்றத்தைச் சொல்கிறது.
அடுத்த காட்சிவாய்பாடு காட்டும் முழுமாற்றத்திற்கு ஏற்ப,
தமிழ் உடல் ஒவ்வொரு $1$ ஐயும் $0$ ஆகவும் ஒவ்வொரு $0$ ஐயும்
$1$ ஆகவும் மாற்றுகிறது.

\textbf{OLSIZ-010.} மாற்றுக் குறைத்தல் நிறுவலில்,
$f(N)$ இன் வெளியீட்டை மீண்டும் $s_k$ என்று உறையவைக்கப்பட்ட ஆங்கில
மூலம் பெயரிடுகிறது; இங்கும் $k$ கட்டுப்படுத்தப்படவில்லை.
உறுப்புவாரி நிபந்தனை $N$ இன் இயல்புக்குறித் தொடரைத் தனித்த முறையில்
தீர்மானிப்பதால், தமிழ் உடல் வெளியீட்டையும் அதன் உறுப்புகளையும்
$s$, $s(n)$ என்ற ஒரே கட்டுப்பட்ட பெயரால் குறிக்கிறது.

\textbf{TA-SIZ-LOCAL-003.} முதன்மை மற்றும் மாற்றுக் குறைத்தல்
பிரிவுகளில் உள்ள இரு வேறுபட்ட பயிற்சிகளுக்கும் உறையவைக்கப்பட்ட ஆங்கில
மூலம் ஒரே முழு அடையாளமான
\texttt{sfr:siz:red:prob:nat-nat} ஐ வழங்குகிறது. இரு பிரிவுகளையும்
ஒரே நூலில் சேர்க்கும்போது இது இரட்டைச் சுட்டியாகிறது.
மாற்றுப் பிரிவின் பெயரிடைவெளிக்கு ஏற்ப, தமிழ் மாற்றுப் பயிற்சியின்
அடையாளத்தை \texttt{sfr:siz:red-alt:prob:nat-nat} எனத் திருத்துகிறது;
முதன்மைப் பயிற்சியின் அடையாளம் மாறவில்லை.

\subfile{../translation/content/sets-functions-relations/arithmetization/arithmetization.tex}
\section*{எண்முறைமைகளின் கட்டமைப்பு பற்றிய பதிப்புக் குறிப்புகள்}
இந்தக் குறிப்புகள் மூல நூலின் பகுதிகள் அல்ல.

\textbf{TA-ARI-001.} விகிதமுறு எண்களின் வரிசையை வரையறுக்கும்
விளக்கத்தில், $r\leq s$ என்பதற்கு $s-r$ எதிர்மமற்றது என்று சரியாகக்
கூறிய உடனே, உறையவைக்கப்பட்ட ஆங்கில மூலம் ``$r-s$ ஐ''
எதிர்மமற்ற மேலெண் கொண்ட பின்னமாக எழுதலாம் என்கிறது.
அடுத்த காட்சிவாய்பாடும் முந்தைய கூற்றும் $s-r$ என்பதையே பயன்படுத்துகின்றன.
தமிழ் உடல் அந்த ஒரே இடத்தை $s-r$ எனத் திருத்துகிறது.

\textbf{TA-ARI-002.} வெட்டுகளின் கணத்தின் முழுமையை நிறுவும்
முதல் குறிப்பில், $S$ இல் குறைந்தது ஒரு வெட்டு இருப்பதற்குக் காரணம்
$S$ க்கு மேல்வரம்பு இருப்பது என்று உறையவைக்கப்பட்ட ஆங்கில மூலம்
கூறுகிறது. ஆனால் அந்த முடிவுக்குத் தேவையானதும் அங்கு முன்கொள்ளப்பட்டதும்
$S$ வெற்றற்றது என்பதே. தமிழ் உடல் காரணத்தை $S$ வெற்றற்றது எனத்
திருத்துகிறது; அடுத்த தனித்த முடிவுக்கு மேல்வரம்பு முன்கோள் அப்படியே உள்ளது.

\textbf{TA-ARI-003.} கோஷி கட்டமைப்பில், சமானத் தொடர்பை
வரையறுத்தபின் மெய்யெண்களை ``சமானத் தொடர்புகளுடன்'' அடையாளப்படுத்துவதாக
ஆங்கில மூலம் கூறுகிறது. அடுத்த வாக்கியமும் முந்தைய இரு கட்டமைப்புகளும்
கோருவது சமான வகுப்புகளே. தமிழ் உடல் ``சமான வகுப்புகளுடன்'' எனத்
திருத்துகிறது.

\textbf{TA-ARI-004.} கோஷி சமான வகுப்பு நேர்மமானது என்பதற்கான
வரையறையில், உறையவைக்கப்பட்ட ஆங்கில மூலம் அதை விகிதமுறு பூச்சியம்
$0_\Rat$ உடன் ஒப்பிடுகிறது. வரையறுக்கப்படுவது மெய்யெண்களின்
சமான வகுப்புகள் என்பதால், தமிழ் உடல் மெய்யெண் பூச்சியம் $0_\Real$
எனத் திருத்துகிறது.

\textbf{TA-ARI-005.} கோஷி கட்டமைப்பின் புலத் தேற்றமும் அதன்
பயிற்சியும், ``கோஷி தொடர்கள்'' ஒரு வரிசைப்படுத்தப்பட்ட புலத்தை
அமைக்கின்றன எனச் சுருக்கமாகக் கூறுகின்றன.
முன்னுள்ள கட்டமைப்பு $\Realequiv$ இன் கீழான சமான வகுப்புகளையே
மெய்யெண்களாக வரையறுக்கிறது; மூலத் தொடர்களை நேரடி சமத்துவத்துடன்
எடுத்தால் புலம் கிடைக்காது. தமிழ் தேற்றமும் பயிற்சியும்
``கோஷி தொடர்களின் சமான வகுப்புகள்'' எனத் தெளிவாக்குகின்றன.
பின்னுள்ள நிறுவல் வரைவில் மூலத்தின் பிரதிநிதி-தொடர் குறுக்குவழி
அப்படியே வைக்கப்பட்டுள்ளது.

\textbf{TA-ARI-006.} வெட்டுகள் மூலம் மெய்யெண்களை அமைக்கும் பிரிவில்,
கழித்தலுக்கும் வகுத்தலுக்கும் கொடுக்கப்பட்ட இரு வரையறைகளும்
உறையவைக்கப்பட்ட மூலத்தில் கருத்துரையாக்கப்பட்டுள்ளன; அவை செயலிலுள்ள
நூல் உரையாக இல்லை. ஆனால் சரிபார்ப்பு இணைப்பு பின்னர்
$\alpha-\beta$, $\alpha\div\beta$ ஆகிய இரண்டும் வெட்டுகள் எனச்
சரிபார்க்குமாறு கூறுகிறது. தமிழ் மூலம் கருத்துரையாக்கப்பட்ட வரையறைகளையும்
பின்னுள்ள குறிப்பையும் அவற்றின் மூல நிலையிலேயே வைக்கிறது;
விடுபட்ட செயல்வரையறையை உருவாக்கிச் சேர்க்கவில்லை.

\subfile{../translation/content/sets-functions-relations/infinite/infinite.tex}
\section*{முடிவிலிக் கணங்கள் பற்றிய பதிப்புக் குறிப்புகள்}
இந்தக் குறிப்புகள் மூல நூலின் பகுதிகள் அல்ல.

\textbf{TA-INF-001.} டெடெகிண்ட் மூடல் பற்றிய துணைத்தேற்றம்,
``எந்தச் சார்பு $f$ க்கும், எந்த $o\in A$ க்கும்'' என்கிறது;
ஆனால் $A$ அங்கு அறிமுகப்படுத்தப்படவில்லை. அடுத்த நிறுவலில்
$\ran{f}\cup\{o\}$ என்பது $f$-மூடப்பட்டதாக இருப்பதற்குத் தேவையான,
பின்னுள்ள பயன்பாட்டிலும் உள்ள, மறைமுக நிபந்தனை $f:A\to A$ என்பதாகும்.
தமிழ் உடல், $f$ இன் சார்பகமும் வீச்சகமும் ஒரே கணம் $A$ எனத்
தெளிவாக்குகிறது; காட்சிவாய்பாடுகளையும் முடிவையும் மாற்றவில்லை.


% The translated propositional part imports proof-system and completeness
% chapters by design. This cumulative reader includes the translated syntax
% here, then imports the proof-system material once in its FOL form below.
\tagfalse{prfSC,prfND,prfTab,prfAX}
\begingroup
\let\olreaderoriginalimport\olimport
\RenewDocumentCommand \olimport { s o m o }{%
  \ifstrequal{#3}{proof-systems}{}{%
    \ifstrequal{#3}{completeness}{}{%
      \IfBooleanTF{#1}{%
        \IfNoValueTF{#2}{%
          \IfNoValueTF{#4}{\olreaderoriginalimport*{#3}}{\olreaderoriginalimport*{#3}[#4]}%
        }{%
          \IfNoValueTF{#4}{\olreaderoriginalimport*[#2]{#3}}{\olreaderoriginalimport*[#2]{#3}[#4]}%
        }%
      }{%
        \IfNoValueTF{#2}{%
          \IfNoValueTF{#4}{\olreaderoriginalimport{#3}}{\olreaderoriginalimport{#3}[#4]}%
        }{%
          \IfNoValueTF{#4}{\olreaderoriginalimport[#2]{#3}}{\olreaderoriginalimport[#2]{#3}[#4]}%
        }%
      }%
    }%
  }%
}
\subfile{../translation/content/propositional-logic/propositional-logic.tex}
\endgroup

\tagtrue{FOL}
\subfile{../translation/content/first-order-logic/proof-systems/proof-systems.tex}
\subfile{../translation/content/first-order-logic/sequent-calculus/sequent-calculus.tex}

% Frege1884's bibliography note cites this translation; BibTeX does not
% discover citations emitted only from a generated .bbl.
\nocite{Frege1953}
\bibliographystyle{../upstream/bib/natbib-oup}
\bibliography{../upstream/bib/open-logic}
\end{document}
